% Spettrometro di massa del test del respiro (ud6-magnetismo-spettrometro): ioni 12C+ e 13C+
% entrano con la stessa velocità nel campo, che entra nel foglio, percorrono mezzo giro e
% arrivano sul rivelatore. Disegno in scala 1:1: 2 r_12 = 6,25 cm, 2 r_13 = 6,75 cm.
\begin{tikzpicture}[>=Stealth, font=\small]
  % regione del campo
  \fill[medica!6] (-7.6,0) rectangle (0.9,3.9);
  \foreach \x in {-7.2,-6.4,...,0.6} \foreach \y in {0.4,1.2,...,3.8}
    \node[medica!60, font=\scriptsize] at (\x,\y) {$\otimes$};
  \node[medica, fill=medica!6, inner sep=1pt, anchor=east] at (0.85,3.6) {$\vec B$ entra};
  % parete con la fenditura d'ingresso, e rivelatore
  \draw[very thick, richiamo] (-7.6,0) -- (-7.05,0);
  \draw[very thick, richiamo] (-6.0,0) -- (-0.08,0);
  \draw[very thick, richiamo] (0.08,0) -- (0.9,0);
  \fill[lampo!70] (-7.05,-0.12) rectangle (-6.0,0.06);
  \node[lampo, anchor=east] at (-7.65,0) {rivelatore};
  % ioni in arrivo
  \draw[->, very thick] (0,-0.8) -- (0,-0.12) node[midway, right] {$\vec v$};
  \fill (0,-0.8) circle (2pt);
  % traiettorie
  \draw[very thick, esempio] (0,0) arc[start angle=0, end angle=180, radius=3.125];
  \draw[very thick, formula] (0,0) arc[start angle=0, end angle=180, radius=3.375];
  \draw[->, very thick, esempio] (-3.125,3.125) arc[start angle=90, end angle=110, radius=3.125];
  \draw[->, very thick, formula] (-3.375,3.375) arc[start angle=90, end angle=110, radius=3.375];
  \node[esempio, fill=medica!6, inner sep=1pt] at (-1.8,1.9) {$^{12}\mathrm{C}^{+}$};
  \node[formula, fill=medica!6, inner sep=1pt] at (-0.3,2.6) {$^{13}\mathrm{C}^{+}$};
  % distanza fra i punti d'arrivo
  \draw[thin, dashed] (-6.75,-0.12) -- (-6.75,-1.65);
  \draw[thin, dashed] (-6.25,-0.12) -- (-6.25,-1.1);
  \draw[thin, dashed] (0,-0.85) -- (0,-1.65);
  \draw[->] (-7.5,-0.45) -- (-6.75,-0.45);
  \draw[->] (-5.5,-0.45) -- (-6.25,-0.45);
  \node[inner sep=1pt, anchor=west] at (-5.45,-0.45) {\SI{5,0}{\milli\metre}};
  % diametri
  \draw[<->, esempio] (-6.25,-1.05) -- (0,-1.05);
  \node[esempio, fill=white, inner sep=1pt] at (-3.1,-1.05) {$2r_{12}$};
  \draw[<->, formula] (-6.75,-1.6) -- (0,-1.6);
  \node[formula, fill=white, inner sep=1pt] at (-3.4,-1.6) {$2r_{13}$};
\end{tikzpicture}
